% RESEARCH_DRAFT_SHA256: fb280486eaefa3cf81c0f0875ccdafaa345cb4e6e8727d55c3ef77c8274a93ae
% Layout and source-reference wrappers expanded for the web renderer.
\documentclass[11pt]{article}
\usepackage[margin=25mm]{geometry}
\usepackage{amsthm}
\usepackage{fontspec,unicode-math}
\setmainfont{Latin Modern Roman}
\setmathfont{Latin Modern Math}
\usepackage{amsmath,amssymb,mathtools,xurl,hyperref}
\hypersetup{colorlinks=true,urlcolor=blue,linkcolor=blue}
\setcounter{secnumdepth}{0}
\setlength{\parindent}{0pt}
\setlength{\parskip}{5pt}
\setlength{\emergencystretch}{3em}
\newcommand{\sref}[2]{\hyperlink{source:#1:#2}{[S#2]}}
\newcommand{\cataloguescope}[1]{\paragraph{Source record.}#1}
\newcommand{\floor}[1]{\left\lfloor#1\right\rfloor}
\newtheorem{proposition}{Proposition}
\newtheorem{lemma}[proposition]{Lemma}
\begin{document}
\paragraph{Research manuscript.} OpenAI \texttt{gpt-6-astra}, reasoning \texttt{ultra}; 11 October 2026.
\paragraph{Preserved source body.} The complete catalogue mathematical text and sources follow with layout and source links expanded, including source qualifications and its historical preparation notes. New research begins at the section headed Research attempt.

\section{A semantics-bound certified bridge between Lean and computer algebra}\label{statement}

\subsection{Definitions and mathematical statement}

\paragraph{Shared notation.}Standard number systems are $\mathbb N,\mathbb Z,\mathbb Q,\mathbb R,\mathbb C$. Specialized notation and index ranges follow the cited source and the definitions below.


A computer algebra system performs symbolic computations on mathematical expressions and algebraic structures. A formal theorem prover checks derivations in an explicit logical framework. A certificate is additional data whose correctness can be checked; connecting a symbolic computation to a proof requires relating its encoded input and result to formal mathematical definitions.
\paragraph{Record scope and qualifications.}
The catalogue title is retained as metadata. Named certificates, contracts, bounds or frameworks in that title are not asserted as results of the original source extract.
The extract does not specify a computation, logic, interface, trust boundary, or verification theorem. It includes a software-design discussion as well as possible implementation work.

Catalogue screening limitation: Workshop planning or software infrastructure agenda; the supplied passage does not identify a determinate mathematical target.

The source snapshot labels this record partially solved; this is catalogue metadata rather than independent verification of a mathematical result or the current completion of a project.
\paragraph{Statement recorded in the supplied source.}
The group has the express goal of finding ways in which current computer algebra systems can interact with formal theorem provers, such as Lean, and vice versa. The guiding motivation for this project is to discuss to what extent Lean should operate as a CAS in it's own right. This groups concrete goals could be anything from the actual development of CAS technology in Lean, or could be to create an argument about the extent to which Lean should be used as a computer algebra system at all.
The primary focus of this group should be to decide to what extent these mathematical technologies should mix and to come up with ways to implement these solutions. \href{http://aimpl.org/cyberinfrastructure/6/}{[S2]}
\paragraph{Source status.}
The supplied catalogue labels this source record partially solved; that metadata was checked against the supplied snapshot on 2026-10-11. The label is preserved for provenance and is not an independent verification of present open status or mathematical correctness.
\subsection{Short English statement}
Explore interactions between computer algebra systems and formal theorem provers, including Lean.
\subsection{Sources}
\begin{itemize}
\item[S1] Ulam AI and the UnsolvedMath Contributors, UnsolvedMath: A Curated Collection of Open Mathematics Problems, catalogue entry. Catalogue curation and metadata: CC BY 4.0. The source extract is retained; definitions, English summaries and qualifications are editorial additions.
\url{https://huggingface.co/datasets/ulamai/UnsolvedMath}.
\item[S2] AIM Workshop Problem Lists, source linked by the catalogue.
\url{http://aimpl.org/cyberinfrastructure/6/}.
\end{itemize}
\paragraph{Preparation and review.}
The original source and UnsolvedMath contributors retain attribution for the source extract. Codex (GPT-6; exact serving revision not exposed) assisted with definitions, English summaries, scope notes and formatting. No new human mathematical authorship or independently verified solution is asserted. Catalogue identity checked against the complete supplied registry and current repository. Distinct upstream records, including repeated questions and broad or damaged extracts, are retained for complete coverage of the source CS category. A repeated record is not a new independent conjecture. Complete mathematical-equivalence and current-literature review has not been performed. Preparation and final review follow contribution prompts C and D, with source gaps disclosed. Structural and rendering checks do not establish a new theorem or open status.
\label{end:statement}

\section{Research attempt}
\subsection{Literature, source context and scope}
The complete catalogue body above retains the original agenda, source status, and historical preparation. The AIMPL problem page could not be retrieved. The supplied complete source transcription was read, and the accessible primary AIM workshop announcement was checked for context. The announcement is not a substitute for verification of the inaccessible problem page. The catalogue's earlier proposed frameworks and research summaries are context, not established solutions credited to this attempt. The following is an explicitly limited investigation of one route; it neither substitutes a formalized toy model for the entire source agenda nor claims empirical or deployed-software results. Elementary auxiliary facts are not new original-target progress.

\subsection{An exact certificate route from a CAS to a checker}
For one sharply specified interoperability task, let $K$ be a field and $R=K[x_1,\ldots,x_n]$ with exact coefficient operations. Given $g_1,\ldots,g_s,f\in R$, an untrusted computer algebra system may propose coefficients $q_1,\ldots,q_s\in R$ certifying that $f$ belongs to the ideal generated by the $g_i$.
\begin{proposition}
Checking the polynomial identity $f=\sum_iq_ig_i$ is a sound certificate of ideal membership. At every common zero $a$ of the $g_i$ in any field extension of $K$, such a certificate also proves $f(a)=0$.
\end{proposition}
\begin{proof}
By definition the generated ideal consists of finite sums $\sum_i r_i g_i$ with $r_i\in R$. The checked coefficients exhibit such a sum. Polynomial evaluation at $a$ respects addition and multiplication, so the identity gives $f(a)=\sum_i q_i(a)g_i(a)=0$. To check the identity algebraically, expand both sides using distributivity, collect coefficients of each monomial, and compare them in $K$. Polynomials are finite coefficient functions on exponent vectors, so agreement of every coefficient is exactly equality in $R$.
\end{proof}
For $R=\mathbb Q[x,y]$, take $g_1=x^2-y$, $g_2=xy-1$, and $f=x^3-1$. The certificate $(q_1,q_2)=(x,1)$ works because
\[
x(x^2-y)+(xy-1)=x^3-1.
\]
This makes one bridge technically concrete, but it cannot simply be reused to certify a different claim. For example, verifying that a returned factor list multiplies to $x^2-1$ does not certify irreducibility: the one-item list $(x^2-1)$ passes the multiplication test even though $(x-1)(x+1)$ is a nontrivial factorization over $\mathbb Q$.
\subsection{Verification and remaining gaps}
A sparse exact polynomial check expanded the certificate and found zero residual coefficients. It also checked the nontrivial factorization counterexample. No Lean code, proof-kernel integration, general CAS protocol, or semantic translation between differing coefficient domains was implemented. Exact ring operations and a matching statement are essential; this is a familiar certificate design, not a proof that Lean should absorb a CAS or that all CAS answers are certified.

\section{References and source attribution}
Primary source: \url{http://aimpl.org/cyberinfrastructure/6/}. Attempted primary AIMPL page retrieval; failed with timeout or gateway error. The supplied full source transcription was read; the inaccessible live page was not independently inspected.

\url{https://aimath.org/pastworkshops/cyberinfrastructure.html}. Read the primary December 2023 workshop announcement; it describes the open-source mathematical infrastructure agenda.

All retrievals were checked on 11 October 2026. Original contributor and catalogue attribution remains in the preserved Sources section. The elementary arguments below are self-contained; their exposition is not a claim of first discovery.


\paragraph{Integration review.} Codex (GPT-6; exact serving revision not exposed) reviewed the written arguments and expanded layout and source-reference wrappers. Repository metadata and the production web renderer were checked using Python 3.13 and JavaScript with MathJax 3.2.2. This model-assisted review is not independent human proof verification; no PDF compilation or formal proof checking is asserted.

\section{Completion Estimate}
\noindent\textbf{COMPLETION: 0\%}
No substantiated new progress toward the entire original target is claimed. The worked results and failed route are delimited above; a known result, illustrative model, or auxiliary obstruction is not counted as original-target completion.

\section{Final status and attribution}
\textbf{Attempt status: unresolved.} This describes this research attempt, not a new ruling on the source catalogue's status or on external literature claims.
The mathematics and this self-audit were generated by OpenAI \texttt{gpt-6-astra}, reasoning \texttt{ultra}, on 11 October 2026. Codex (GPT-6; exact serving revision not exposed) orchestrated the batch and contributed formatting/repository review. No human mathematical authorship, independent human proof checking, or formal verification is claimed. The historical preparation paragraphs in the reproduced catalogue body concern that body only.

Contribution prompt A was applied to the complete definition and supplied source record, followed by the adversarial checks of prompt B. Local reproduction files are \texttt{batch013/checks.py}, \texttt{computations.json}, and the per-ID audit JSON. They were not public links at preparation time. Numerical checks are finite and do not prove unrestricted results. No LaTeX engine was found locally; compilation was not performed by this research agent.

\end{document}
